\documentclass[10pt]{article}

% Specify the margins and text size.
\setlength{\textwidth}{6.4in}
\setlength{\textheight}{9.5in}
\setlength{\oddsidemargin}{0pt}
\setlength{\evensidemargin}{0pt}
\setlength{\topmargin}{0pt}
\setlength{\hoffset}{.05in}
\setlength{\voffset}{-1in}

\setlength{\parskip}{5pt}
\setlength{\parindent}{0pt}

% Load some fonts and symbol packages
\usepackage{latexsym}
\usepackage{pifont}       % contains 'star' symbol for counterinsurgency handbook title
\usepackage{yfonts} 
\usepackage{amsmath}
\usepackage{amsfonts}

\usepackage{graphicx}     % actually, this is loaded by pstricks
\usepackage[T1]{fontenc}
\usepackage{ifthen}
\usepackage{pstricks,pst-grad,pst-text,pst-node,multido,pst-plot,calc,pst-3dplot}
%\usepackage[all]{xy}
%\usepackage{animate}

% The hyperref package inserts links into the pdf.
\definecolor{MyLinkColor}{rgb}{.1,.2,1}
\definecolor{MyCiteColor}{rgb}{.1,1,.2}
\definecolor{MyURLColor}{rgb}{.4,.4,.4}
\usepackage[backref=true,pagebackref=false,hyperindex,colorlinks=true,
  linkcolor=MyLinkColor,urlcolor=MyURLColor]{hyperref}


% The tweaklist package is something I found on the web.  It provides a simple interface
% for making changes to spacing used in the itemize and enumerate environments.  Comment
% this out if you don't care to use tweaklist.
\usepackage{tweaklist}
\renewcommand{\itemhook}{\setlength{\parskip}{2pt}\setlength{\parsep}%
{1pt}\setlength{\topsep}{0pt}\setlength{\itemsep}{0pt}}

\newcommand{\U}{\underline{\hspace{5pt}}}

\usepackage{listings}
\newcommand{\Z}{\hphantom{0}}

\begin{document}
\pagestyle{empty}
\lstset{language=R, showspaces=false, showstringspaces=false}

\href{http://www.su.edu}{\includegraphics[height=1.75cm]{sulogo.eps}}
\vspace{-1.79cm}

{{\ }\hfill\small
\begin{tabular}{cl}
& Math 207\\
& Introduction to Statistics\\
\end{tabular}
}
\setlength{\baselineskip}{1.05\baselineskip}
\medskip

\begin{center}
\textbf{\large   The Normal Approximation for Probability Histograms}
\end{center}
\smallskip

\newcommand{\HH}{\hspace{20pt}\hphantom{a) } }

1. Five coins are to be tossed and the number of heads counted.  
The number of heads is like 
the sum of five draws (with replacement) from a box with two tickets $1$ 
(for heads) and $0$
(for tails).

\HH \lstinline!> source("http://www.adjoint-functors.net/su/web/314/R/BoxSimulations.R");!

\hspace{20pt} a) Generate data like that in Table~1 on page 309 of your text.

\HH \lstinline!> table(boxSimulation(c(0, 1), 5, 100))!

Are the numbers generated by R the same as those shown in the table?  Why or why not?
Enter the command again.  Does R generate the same numbers? Why or why not?

\hspace{20pt} b) Make a histogram of the data from 100 repetitions of tossing five coins.

\HH \lstinline!> boxHistogram(c(0,1), 5, 100, breaks=-0.5+(0:6))!

\hspace{20pt} c) Increase the number of repetitions from 100 to 1000.  

\HH \lstinline!> boxHistogram(c(0,1), 5, 1000, breaks=-0.5+(0:6))!

How many repetitions does it take before the histograms no longer appear to change?

\hspace{20pt} d) Increase the number of coins from 5 to 10.

\HH \lstinline!> boxHistogram(c(0,1), 10, 1000, breaks=-0.5+(0:11))!

As you increase the number of coins, keeping the number of repetitions fixed, does the 
histogram tend to look more like a normal curve or less?
\vspace{.5in}

2. Two dice are to be rolled and the sum of the number of spots counted.  

\HH \lstinline!> source("http://www.adjoint-functors.net/su/web/314/R/diceSimulations.R");!

\hspace{20pt} a) Generate data like that in Table~2 on page 310 of your text.  

\HH \lstinline!> diceSimulation(100)!

Are the numbers generated by R the same as those shown in the table?  Why or why not?
Enter the command again.  Does R generate the same numbers? Why or why not?

\hspace{20pt} b) Plot a histogram of the data from 100 repetitions of rolling two dice.

\HH \lstinline!> diceHistogram(100)!

Does your plot differ from the one at the top of Figure~1 on page 311 of your text?  Increase the
number of repetitions to try to reproduce the other plots in Figure~1.  For example:

\HH \lstinline!> diceHistogram(1000)!

How many repetitions does it take before the histograms no longer appear to change?

\hspace{20pt} c) Change the color of the histogram bars then try to discover some other colors that R knows.

\HH \lstinline!> diceHistogram(1000, color="green")!

\hspace{20pt} d) Change the number of dice.

\HH \lstinline!> diceHistogram(1000, dice=3)!

As you increase the number of dice, keeping the number of repetitions fixed, does the histogram tend to
look more like a normal curve or less?


\vfill
\eject          
{\ }

3. Two dice are to be rolled and the product of the number of spots counted.  

\hspace{20pt} a) Generate data for 100 repetitions of the experiment.

\HH \lstinline!> diceSimulation(100, diceSum=FALSE)!

Enter the command again.  Does R generate the same numbers? Why or why not?

\hspace{20pt} b) Plot a histogram of the data from 100 repetitions of rolling two dice.

\HH \lstinline!> diceHistogram(100, diceSum=FALSE)!

Does your plot differ from the one at the top of Figure~2 on page 313 of your text? 
Why are there gaps in the histogram?

 Increase the
number of repetitions to try to reproduce the other plots in Figure~2.  For example:

\HH \lstinline!> diceHistogram(1000, diceSum=FALSE)!

How many repetitions does it take before the histograms no longer appear to change?
Do the histograms look like a normal curve?  Does your answer contradict The 
Central Limit 
Theorem (page 325)?
\vspace{.5in}

4. Consider sums of draws from the box describe in Figure~9 on page 322 of your text.

\hspace{20pt} a) Plot a histogram of data from 100 repetitions of sums of 25 draws.

\HH \lstinline!boxHistogram(c(1,2,9), 25, 100, ymax=0.1, breaks=160)!

Does the histogram look like  a normal curve?

\hspace{20pt} b) Find formulas for the expected value, maximum sum, minimum sum and
standard error for sums of $n$ draws.

\hspace{20pt} c)  Increase the number of repetitions to get a plot that looks like the 
first one in Figure~9.  You may need to decrease ymax and increase breaks.

\hspace{20pt} d)  Repeat c) for sums of 50 and 100 draws as shown in Figure~9.
\medskip

5. Do exercise 5 from page 314 of your text.  use \lstinline!boxHistogram! to help and to 
check your answers.
\medskip

6. Do exercise 5 from page 324 of your text.  use \lstinline!boxHistogram! to help and to 
check your answers.
\medskip

7. Do exercise 6 from page 324 of your text.  use \lstinline!boxHistogram! to help and to 
check your answers.

\end{document}

